\begin{table}[t]
\centering
\caption{Component-level inference latency profile measured on standard clinical workstation hardware.}
\label{tab:runtime_latency}
\small
\begin{adjustbox}{max width=\linewidth}
\begin{tabular}{lllcc}
\toprule
\textbf{Pipeline Component} & \textbf{Target Operation} & \textbf{Input Modality} & \textbf{Mean Latency (ms)} & \textbf{95th Pct (ms)} \\
\midrule
TwinStager (Model A)        & Multi-tier syndromic staging      & Demographics, Cognition, Fluid, MRI & 36.79 ms & 42.10 ms \\
Progression Engine (Model B)& Multi-horizon forecast (6/12/24m) & 10 core clinical anchors            & 58.20 ms & 64.50 ms \\
Lifestyle Modifier          & Relative-risk multiplier slider   & 10 lifestyle \& health history flags& 0.072 ms & 0.095 ms \\
TreeSHAP Stage Drivers      & Local feature attribution         & Active CatBoost tree leaves         & 25.48 ms & 29.80 ms \\
Scenario Lab Simulation     & Complete counterfactual roll-out  & Baseline twin + parameter overrides & 222.83 ms & 245.00 ms \\
Full Twin Initialization    & End-to-end twin construction      & Complete patient intake payload     & 343.37 ms & 382.10 ms \\
\midrule
\multicolumn{5}{l}{\textit{Benchmark Environment: Intel Core i5-10300H CPU @ 2.50 GHz, 16 GB RAM, Windows 10 x64, Python 3.11, CatBoost 1.2.}} \\
\bottomrule
\end{tabular}
\end{adjustbox}
\end{table}
